\begin{abstract}

Most uncertainty research in medical image segmentation stops at
\emph{estimation}---producing variance maps and reporting calibration
metrics---without specifying what a clinician should \emph{do} with them.
The missing layer is \emph{decision}: a policy that converts uncertainty
into action. \textbf{Uncertainty is not useful unless it informs action.}

We recast uncertainty-aware segmentation as a two-stage pipeline---estimation
followed by decision---and show that the stages interact strongly enough
that optimizing either alone is insufficient. On retinal vessel segmentation
(DRIVE, STARE, CHASE\_DB1), we compare Monte Carlo Dropout and Test-Time
Augmentation (TTA) as uncertainty sources, each paired with three deferral
policies of increasing sophistication: global thresholding, image-adaptive
thresholding, and a confidence-aware rule that weighs uncertainty by
prediction borderline-ness.

The right pairing removes roughly 80\% of segmentation errors at only
25\% pixel deferral; a confidence-aware variant removes 55\% of errors at
just 12\% deferral. TTA achieves Unc-AUROC 0.881 versus 0.722 for MC
Dropout and runs $3.1\times$ faster. Temperature scaling---the default
calibration fix---does not improve downstream decisions, and uncertainty
quality transfers robustly under domain shift even as Dice drops by 6.5
points. These results establish a concrete evaluation principle:
\textbf{uncertainty methods should be judged by the decisions they enable,
not by how well their probabilities are calibrated in isolation.}

\end{abstract}
